% ===========================================================================
\section{The binary erasure channel}\label{sec:bec}
\warmth{0}\quad\whisper{Losing a bit is cheaper than being lied to about it.}

\begin{strand}
The erasure channel loses a fraction $\alpha$ of the bits but never corrupts
one: the receiver always knows \emph{which} bits were lost. Its capacity is
$1-\alpha$ --- exactly the fraction of channel uses that survive. Compare this
with the BSC, where a flip probability of $\alpha$ costs $h(\alpha)$, which is
much more.
\end{strand}

\trigger{When one output symbol means ``no information'', condition on the event that it occurred.}

\block{The channel}

\begin{definition}[binary erasure channel]\label{def:bec}
The \keyword{binary erasure channel} $\mathrm{BEC}(\alpha)$ has $\X=\{0,1\}$ and
$\Y=\{0,1,e\}$, where $e$ is the \keyword{erasure} symbol, and
\[
  p(0\mid0)=p(1\mid1)=1-\alpha,
  \qquad
  p(e\mid0)=p(e\mid1)=\alpha .
\]
\begin{center}
\begin{tikzpicture}[x=2.8cm,y=1.2cm,>=latex,font=\small]
  \node (x0) at (0,2) {$0$};
  \node (x1) at (0,0) {$1$};
  \node (y0) at (1,2) {$0$};
  \node (ye) at (1,1) {$e$};
  \node (y1) at (1,0) {$1$};
  \draw[->,indigo] (x0) -- node[above,font=\scriptsize]{$1-\alpha$} (y0);
  \draw[->,indigo] (x1) -- node[below,font=\scriptsize]{$1-\alpha$} (y1);
  \draw[->,madder] (x0) -- node[pos=0.6,above,sloped,font=\scriptsize]{$\alpha$} (ye);
  \draw[->,madder] (x1) -- node[pos=0.6,below,sloped,font=\scriptsize]{$\alpha$} (ye);
\end{tikzpicture}
\end{center}
No bit is ever complemented, so an output in $\{0,1\}$ identifies the input with
certainty; an output $e$ says nothing at all. Every row is a permutation of
$(1-\alpha,\alpha,0)$, hence $\Ent(Y\mid X)=\answerin[1.6cm]{h(\alpha)}$.
\studentrecall{Each row is $(1-\alpha,\alpha)$ on two symbols, so the row entropy is $h(\alpha)$.}
\end{definition}

\block{Capacity}

\begin{theorem}[capacity of the BEC]\label{thm:bec-capacity}
The capacity of $\mathrm{BEC}(\alpha)$ is
\[
  C=1-\alpha\quad\text{bits per channel use},
\]
attained by the uniform input distribution.
\end{theorem}

\begin{proof}
Write, as always,
\[
  C=\max_{p(x)}\MI(X;Y)
   =\max_{p(x)}\bigl[\Ent(Y)-\Ent(Y\mid X)\bigr]
   =\max_{p(x)}\Ent(Y)-h(\alpha).
\]
To compute $\Ent(Y)$, let $E=\mathbf{1}\{Y=e\}$ be the indicator of an erasure.
Since $E$ is a function of $Y$, we have $\Ent(Y,E)=\Ent(Y)$, and the chain rule
gives
\[
  \Ent(Y)=\Ent(Y,E)=\Ent(E)+\Ent(Y\mid E).
\]
Now $E\sim\operatorname{Bernoulli}(\alpha)$, so $\Ent(E)=h(\alpha)$. Given
$E=1$ the output is $e$ with certainty, contributing nothing; given $E=0$ the
output equals the input. Writing $\PP(X=1)=\pi$,
\[
  \Ent(Y\mid E)=\alpha\cdot0+(1-\alpha)\,h(\pi),
  \qquad\text{so}\qquad
  \Ent(Y)=h(\alpha)+(1-\alpha)h(\pi).
\]
Therefore
\[
  C=\max_{\pi}\bigl[h(\alpha)+(1-\alpha)h(\pi)\bigr]-h(\alpha)
   =\max_{\pi}(1-\alpha)h(\pi)
   =1-\alpha,
\]
the capacity being achieved at $\pi=\tfrac12$.
\end{proof}

\begin{remark}
The interpretation is exact rather than approximate: a fraction $\alpha$ of the
transmitted bits is lost, so at most a fraction $1-\alpha$ can get through, and
that much is achievable. Note also that here the two ``obvious'' bounds behave
differently: $\Ent(Y)\le\log3$ is \emph{not} tight, because the erasure symbol
can never be reached more often than $\alpha$ of the time. Splitting $\Ent(Y)$
via the erasure indicator is what replaces the crude bound.
\end{remark}

\block{BSC versus BEC at the same noise level}

\noindent\begin{tabularx}{\linewidth}{@{}l l L@{}}
\toprule
{\color{indigo}channel} & {\color{indigo}capacity} & \multicolumn{1}{l}{\color{madder}what the receiver knows}\\
\midrule
$\mathrm{BSC}(p)$ & $1-h(p)$ & every received bit may be a lie; the location of errors is unknown\\
$\mathrm{BEC}(\alpha)$ & $1-\alpha$ & received bits are certainly correct; erasures are flagged\\
\bottomrule
\end{tabularx}

At $p=\alpha=0.1$ these give $C\approx0.531$ and $C=0.9$: knowing \emph{where}
the damage is worth a lot.

\begin{yourturn}
The key trick was $\Ent(Y)=\Ent(E)+\Ent(Y\mid E)$. State the two facts that make
this identity legitimate and useful here.
\studentrecall{$E$ is a function of $Y$, so $\Ent(Y,E)=\Ent(Y)$; and given $E$ the output is either fixed at $e$ or equal to $X$.}
\ifshowanswers\else\workspace[3]\fi
\answerprose{First, $E=\mathbf{1}\{Y=e\}$ is a deterministic function of $Y$,
so $\Ent(Y,E)=\Ent(Y)$ and the chain rule may be applied to the pair. Second,
conditioning on $E$ splits the output into a degenerate case ($Y=e$, entropy
$0$) and a noiseless case ($Y=X$, entropy $h(\pi)$), which is why
$\Ent(Y\mid E)=(1-\alpha)h(\pi)$.}
\end{yourturn}

\begin{yourturn}
Complete the two capacities from memory:
$\mathrm{BSC}(p)$ has $C=\answerin[2.2cm]{1-h(p)}$ and $\mathrm{BEC}(\alpha)$
has $C=\answerin[2.0cm]{1-\alpha}$.
\studentrecall{$1-h(p)$ for the BSC; $1-\alpha$ for the BEC.}
\end{yourturn}

\recall{Why is $\Ent(Y)\le\log3$ useless for the BEC?}
